\section{Introduction}

As quantum computing hardware continues to advance, so too does its software.
While quantum circuits remain the de facto standard language for low-level quantum programming, they are cumbersome when designing large-scale programs with tens of thousands of gates.
Higher-level abstractions of quantum circuits are on the horizon. 
Diagrammatic reasoning~\cite{Coecke2011} is a natural fit for reasoning about such abstractions since quantum circuits are diagrams. 
Diagrammatic reasoning has already been applied to a variety of circuit-related tasks, including optimization~\cite{Kissinger2020}, simulation~\cite{Sutcliffe2024,Sutcliffe2025}, equivalence checking~\cite{Peham2022}, and error correction~\cite{Khesin2025,Kissinger2024,Rodatz2024}.
Equational theories are collections of equations that serve as toolboxes for performing diagrammatic reasoning.

Control is one such ubiquitous abstraction in quantum algorithm design. Controlled subcircuits appear in essentially every powerful quantum algorithm~(e.g.,~\cite{Kitaev1995,Shor1997,Grover1996,Harrow2009}). 
We generalize the existing notion of control by allowing a subcircuit to be controlled not just by one qubit in a particular basis, but by an arbitrary basis on any number of qubits.
This notion of control is captured by a new binary operation $( \odot )$ called the \textit{control product} which joins two subcircuits together to create a generalized controlled operation.
The simplest example is the $\textsf{CNOT}$, which is created by joining the Pauli $Z$ and $X$ gates together as $Z \odot X = \textsf{CNOT}$.
However unlike the normal notion of control, the control product allows any two subcircuits $U$ and $V$ of arbitrary sizes to control each other as $U \odot V$.
Examples of how such a product can aid in circuit design and verification are found in~\cite{Gidney2017} and~\cite{Schober2024}, respectively.


We present the control product formally in the language of symmetric monoidal categories, as part of a universal algebraic (see~\cite{BaaderNipkow1998}) extension of a prop category~\cite{MacLane1965,Bergholm2011,Clement2023}. We then provide a sound and complete equational theory for performing diagrammatic reasoning on the category $\CatHQC$ of \emph{hierarchical quantum circuits}, that is, quantum circuits with generalized control and circuit exponentiation.
We note that the definition of circuit exponentiation used in this paper is consistent with the power modifier in OpenQASM~3~\cite{Cross2022}.

\vspace{1em}
\noindent
\textbf{Related work}.
The first complete equational theory for uncontrolled quantum circuits was presented in~\cite{Clement2023} and improved in~\cite{Clement2024Extensions,Clement2024Minimal}. 
Two recent works have provided equational theories for standard controlled quantum circuits using controlled props~\cite{Delorme2025} and rig categories~\cite{Heunen2025}. 
Another recent work introduced a related methodology for constructing controlled unitaries via subspace selection and phasing~\cite{Heunen2026}.
The control product we introduce here can be seen as a way to efficiently package exponentially many such selection-and-phasing statements together into one circuit object.
Controlled operations have also recently been studied in the context of quantum control flow~\cite{Yuan2024,Valiron2022} for its applications to quantum computer architectures with a quantum program counter. 

\vspace{1em}
\noindent
\textbf{Structure and contents}.
All proofs are found in the appendices.
\cref{Sect:Background} provides the necessary mathematical background.
\cref{Sect:Semantics} introduces the semantics of control and exponentiation on unitary matrices and provides a characterization in terms of control functors.
\cref{Sect:Hierarchical} constructs the hierarchical quantum circuit language $\CatHQC$.
We describe the inductive construction that is used to generate $\CatHQC$ via closures over an increasing family of nested sub-languages, provide its semantics, and present a sound equational theory.
In \cref{Sect:Completeness} we prove completeness of the equational theory for $\CatHQC$ via a reduction to a known complete equational theory for controlled quantum circuits found in~\cite{Delorme2025}. 
This reduction includes a transpilation procedure which transforms hierarchical circuits into quantum circuits with standard controls.
